\section{Laboratório Prático e Conclusão}

\begin{frame}{Resumo dos Pilares Aprendidos}
    Nesta aula introdutória de 90 minutos, dominamos os quatro alicerces fundamentais da computação gráfica 2D:

    \vspace{0.3cm}
    \begin{columns}
        \begin{column}{0.5\textwidth}
            \begin{itemize}
                \item \textbf{Game Loop:} A espinha dorsal contínua que processa entradas, física e renderização.
                \item \textbf{Plano Cartesiano 2D:} Origem $(0, 0)$ no topo esquerdo e coordenada $Y$ invertida para baixo.
            \end{itemize}
        \end{column}
        \begin{column}{0.5\textwidth}
            \begin{itemize}
                \item \textbf{Controle de Tempo:} Uso de \texttt{Clock.tick(60)} para garantir taxa estável em qualquer hardware.
                \item \textbf{Colisões com Rect:} Otimização e praticidade através do método \texttt{colliderect()}.
            \end{itemize}
        \end{column}
    \end{columns}

    \vspace{0.4cm}
    \begin{block}{Boas Práticas de Engenharia de Jogos}
        \begin{itemize}
            \item Sempre limpe a tela com \texttt{tela.fill()} antes de desenhar.
            \item Trate o evento \texttt{pygame.QUIT} para não congelar o sistema operacional.
            \item Separe a lógica de atualização da renderização visual.
        \end{itemize}
    \end{block}
\end{frame}

\begin{frame}{Laboratório Prático: Desafios para a Turma}
    \begin{block}{Desafio 1: A Bola Multicolorida (Nível Fácil)}
        \begin{itemize}
            \item Execute o arquivo \texttt{desafio1\_bola\_quicando.py}.
            \item Faça a bola quicar continuamente nas 4 paredes da janela.
            \item \textbf{Bônus:} A cada batida, sorteie uma nova cor aleatória com \texttt{random.randint(50, 255)}.
        \end{itemize}
    \end{block}

    \vspace{0.2cm}
    \begin{block}{Desafio 2: Esquiva de Meteoros (Nível Intermediário)}
        \begin{itemize}
            \item Execute o arquivo \texttt{desafio2\_esquiva\_obstaculo.py}.
            \item Mova seu personagem na horizontal para desviar de blocos caindo do topo da tela.
            \item Utilize \texttt{colliderect()} para descontar vidas do jogador se for atingido!
        \end{itemize}
    \end{block}
\end{frame}

\begin{frame}{Próximos Passos no Curso}
    \begin{center}
        \Large{\textbf{O que exploraremos nas próximas aulas?}}
    \end{center}

    \vspace{0.3cm}
    \begin{itemize}
        \item \textbf{Sprites e Grupos:} Uso da classe \texttt{pygame.sprite.Sprite} e \texttt{pygame.sprite.Group} para gerenciar dezenas de entidades com POO.
        \item \textbf{Carregamento de Imagens:} Substituir formas geométricas por Spritesheets e animações de personagens.
        \item \textbf{Sistemas de Partículas e Câmera:} Efeitos visuais dinâmicos e mundos maiores que a tela.
    \end{itemize}

    \vspace{0.5cm}
    \begin{center}
        \textbf{Obrigado!} \\
        \footnotesize{IFCE Campus Tauá -- Curso de Análise e Desenvolvimento de Sistemas}
    \end{center}
\end{frame}
